\documentclass{plantillaPPT}
\titulo{4}{Integral Definida}
\begin{document}
\primeraDiapo
\begin{frame}{Integrales Definidas e Impropias}
\framesubtitle{Práctica 4}
1. Calcule las siguientes integrales definidas

\[(a)\int_0^1 \sqrt{x^2+1} \quad (b)\int_{-2}^{-1} \frac{1}{x(x-1)^2}\,dx \quad (c) \int_1^e \sin(\ln(x)) \,dx\]

\end{frame}
\begin{frame}{Problema 2}
\framesubtitle{Práctica 4}

\((a)\) Sea \(g:\mathbb{R}\to\mathbb{R}\) una función continua tal que
\[\int_{1+x}^{(2x+1)^2} g(t)\,dt = x^4-3x-2\]

Calcular, si es posible, \(g(1)\).
\end{frame}
\begin{frame}{Problema 2}
\framesubtitle{Práctica 4}
\((b)\) Decida si la función \(f:\mathbb{R}\to\mathbb{R}\) definida por
\[f(x) = 
\begin{cases}\displaystyle
    \frac{1}{x-1}\int_0^{x-1} \sin(t^2)\,dt\quad &\text{si } x\neq 1\\
    0 \quad &\text{si } x=1
\end{cases}\]
es derivable y continua en \(X=1\).
\end{frame}

\begin{frame}{Integral Impropia de Tipo I}
    \framesubtitle{Definición}
    \vspace{-0.5cm}

\begin{definicion}
Dada una función \( f:[a, \infty) \to \mathbb{R} \) continua. Se define:
\vspace{0.3cm}
\[
\int_a^{\infty} f(x) \, dx = \lim_{b \to \infty} \int_a^b f(x) \, dx
\]\vspace{0.3cm}

Si el límite existe y es finito, se dice que la integral \textbf{converge}.
\end{definicion}

\end{frame}

\begin{frame}{Integral Impropia de Tipo II}
    \framesubtitle{Definición}

\begin{definicion}
Dada \( f \) continua en \( (a, b) \) excepto en un punto \( c \in (a, b) \). Se define:
\vspace{0.3cm}
\[
\int_a^b f(x) \, dx = \lim_{x \to c^-} \int_a^x f(t) \, dt + \lim_{x \to c^+} \int_x^b f(t) \, dt
\]
\vspace{0.3cm}

Si ambos límites existen y son finitos, se dice que la integral \textbf{converge}.
\end{definicion}

\end{frame}

\begin{frame}{Problema 3}
\framesubtitle{Práctica 4}
3. Sea \(f\) una función derivable en \(\mathbb{R}\) tal que \(f(2)=0\) y \(f(8)=1\). Calcule
\[
\int_1^{\sqrt{3}} \frac{x \cdot f'(3x^2-1)}{1+\left[f(3x^2-1)\right]^2}\,dx
\]
\end{frame}
\begin{frame}{Integrales Impropias}
\framesubtitle{Práctica 4}
4. Calcule, si existe, las siguientes integrales impropias
\[(a)\int_0^1 \frac{\ln(x)}{\sqrt{x}}\,dx \qquad (b)\int_3^{+\infty} \frac{x^2}{x^3-1}\,dx\]
    
\end{frame}


\end{document}
